\documentclass[10pt]{article}
\usepackage[T1]{fontenc}
\usepackage[utf8]{inputenc}
\usepackage[french]{babel}
\usepackage{amsmath,amssymb,mathrsfs,mathpazo}
\usepackage{geometry}
\geometry{a4paper,margin=2cm}
\usepackage{graphicx,array,tabularx,tabularray,colortbl,multirow,multicol,booktabs,enumitem,pifont,ulem,fancybox,diagbox,cancel}
\usepackage{pstricks,pstricks-add,pst-plot,pst-tree,pst-node,pst-text,pst-func,pst-math,pst-bspline,pst-3dplot,pst-eucl,pst-solides3d}
\let\pscalculate\undefined
\usepackage{pst-bezier}
\usepackage{tikz,pgfplots,tkz-tab}
\usetikzlibrary{arrows,arrows.meta,patterns,calc,positioning,trees,decorations.markings,shapes.geometric,shapes.misc,shapes.arrows,intersections,decorations.pathmorphing,decorations.pathreplacing,backgrounds,shadows,babel}
\pgfplotsset{compat=1.18}
\usepgfplotslibrary{fillbetween}
\usepackage[np]{numprint}
\usepackage{listings,makecell,pas-tableur,twoopt}
\let\barre\undefined
\usepackage[pstricks]{tablvar}
\definecolor{darkgreen}{RGB}{0,100,0}
\providecommand{\red}{\color{red}}
\providecommand{\blue}{\color{blue}}
\providecommand{\green}{\color{green}}
\providecommand{\black}{\color{black}}
\providecommand{\white}{\color{white}}
\providecommand{\gray}{\color{gray}}
\providecommand{\cyan}{\color{cyan}}
\providecommand{\magenta}{\color{magenta}}
\providecommand{\yellow}{\color{yellow}}
\providecommand{\violet}{\color{violet}}
\providecommand{\orange}{\color{orange}}
\providecommand{\euro}{\texteuro}
\providecommand{\EUR}{\texteuro}
\providecommand{\ieme}{\textsuperscript{e}}
\providecommand{\degres}{\ensuremath{^\circ}}
\providecommand{\degree}{\ensuremath{^\circ}}
\providecommand{\vv}[1]{\overrightarrow{#1}}
\providecommand{\Coord}[1]{#1}
\providecommand{\decosix}{}
\providecommand{\starredbullet}{\ensuremath{\bullet}}
\pagestyle{empty}
\setlength{\parindent}{0pt}
\usepackage{adjustbox}\DeclareUnicodeCharacter{2236}{:}\DeclareUnicodeCharacter{2190}{\ensuremath{\leftarrow}}
\providecommand{\R}{\mathbb{R}}
\providecommand{\N}{\mathbb{N}}
\providecommand{\Z}{\mathbb{Z}}
\providecommand{\Q}{\mathbb{Q}}
\providecommand{\D}{\mathbb{D}}
\providecommand{\C}{\mathbb{C}}
\providecommand{\e}{{\mathrm{e}}}
\providecommand{\vect}[1]{\overrightarrow{#1}}
\providecommand{\vectt}[1]{\overrightarrow{\mathrm{#1}}}
\providecommand{\vvt}[1]{\overrightarrow{\mathrm{#1}}}
\providecommand{\np}[2][]{{#2~#1}}
\providecommand{\vv}[1]{{\overrightarrow{#1}}}
\renewcommand{\d}{{\mathrm{d}}}
\providecommand{\degres}{{^\circ}}
\providecommand{\degre}{{^\circ}}
\providecommand{\blue}{\color{blue}}
\providecommand{\red}{\color{red}}
\providecommand{\enskip}{\quad}
\providecommand{\llbracket}{\lbrack\!\lbrack}
\providecommand{\rrbracket}{\rbrack\!\rbrack}
\providecommand{\euro}{{\text{€}}}
\providecommand{\no}{{\text{n°}}}
\providecommand{\interval}[2]{{[#1;#2]}}
\providecommand{\intervFO}[2]{{[#1;#2[}}
\providecommand{\dotfill}{\ldots}

\providecommand{\ds}{\displaystyle}
\providecommand{\textding}[1]{#1}
\providecommand{\Oij}{\(\mathrm{O};\vec{\imath},\vec{\jmath}\)}
\providecommand{\Oijk}{\(\mathrm{O};\vec{\imath},\vec{\jmath},\vec{k}\)}
\providecommand{\Ouv}{\(\mathrm{O};\vec{u},\vec{v}\)}
\providecommand{\intervFF}[2]{\ensuremath{[#1;#2]}}
\renewcommand{\euro}{\ifmmode\text{\texteuro}\else\texteuro\fi}
\DeclareUnicodeCharacter{00B0}{\ifmmode^{\circ}\else\textdegree\fi}
\DeclareUnicodeCharacter{20AC}{\ifmmode\text{\texteuro}\else\texteuro\fi}
% Énoncés : Éducation nationale. Transcription et corrigés : APMEP, Denis Vergès et contributeurs. Découpage et classement : CoopMaths, Sébastien Lozano, Éric Elter et contributeurs. Adaptation de la présentation : Maths Explorer.
% CC BY-SA : https://coopmaths.fr/www/about/
% https://coopmaths.fr/alea/static/eam/2026/tex/eam_technologique_2026_06_antilles_2_cor.tex
\begin{document}
\noindent\textit{APMEP / CoopMaths. CC BY-SA. Présentation adaptée par Maths Explorer.}\par\medskip
Puisqu'on interroge une personne au hasard, on est en situation d'équiprobabilité, et les proportions sont donc assimilables à des probabilités.

	\begin{enumerate}
		\item Il y a au total 100 personnes, dont 35 ne vivant pas en couple et n'ayant pas d'enfant,

 		donc \mbox{$P\big(\overline{C}\cap\overline{E}\big) = \dfrac{35}{100} = \np{0.35}$.}

		\item \mbox{$P\big(C\cap\overline{E}\big) = \dfrac{12}{100} = \np{0.12}$} : \quad la probabilité de choisir au hasard une personne vivant en couple et n'ayant pas d'enfant est \np{0.12}.

		\item On cherche \mbox{$P_{E}\big(\overline{C}\big)$}.

		Par définition des probabilités conditionnelles, \mbox{$P_{E}\big(\overline{C}\big) = \dfrac{P\big(E\cap\overline{C}\big)}{P\big(E\big)}$}.

		53 personnes sur les 100 étudiées ont au moins un enfant, donc \mbox{$P\big(E\big) = \np{0.53}$}.

		Parmi les personnes interrogées, 35 ne vivent pas en couple mais ont au moins un enfant, donc \mbox{$P\big(E\cap\overline{C}\big) = \np{0.35}$}

		La probabilité que la personne ne vive pas en couple sachant qu'elle a au moins un enfant est donc \mbox{$P_{E}\big(\overline{C}\big) = \dfrac{35}{53}$.}

		NB : on peut retrouver plus simplement ce résultat en observant que, parmi les 53 personnes ayant au moins un enfant, 35 ne vivent pas en couple.

		\item \begin{enumerate}
			\item Un logement peut donc avoir 1, 2, 3, 4 ou 5 pièces, et les probabilités sont les pourcentages respectifs.

			On a donc la loi de probabilité de $X$ :

			\begin{center}
			\begin{tblr}{colspec={X[l,2cm]*{5}{X[c,1.5cm]}},hlines, vlines}
				$x_i$        & $1$  & $2$ & $3$ & $4$  & $5$  \\
				$P(X = x_i)$ & 0,07 & 0,2 & 0,4 & 0,25 & 0,08\\
			\end{tblr}
			\end{center}

			\item \mbox{$P(X\geqslant 2)$} est l'évènement contraire de \mbox{$P(X = 1)$}, donc \quad \mbox{$P(X\geqslant 2) = 1-\np{0.07} = \np{0.93}$.}

			cela signifie que la probabilité qu'un logement ait un nombre de pièces supérieur ou égal à 2 est \np{0.93}.

			\item L'espérance de $X$ est égale à \quad $\ds \mathrm{E}(X) = \sum_{i = 1}^5 P(X = x_i)\times x_i$.

			$\ds \mathrm{E}(X) = 0,07 \times 1 + 0,2 \times 2 + 0,4 \times 3 + 0,25 \times 4 + 0,08 \times 5 = 0,07 + 0,4 + 1,2 + 1 + 0,4 = 0,47 + 2,2 + 0,4 = 0,47 + 2,6 = 3,07$.

			donc \mbox{$E(X) = \np{3.07}$}.

			Cela signifie que les logements étudiés ont en moyenne 3,07 pièces.

			\emph{Remarque :} Ici, le 3,07 s'interprète comme une moyenne, il n'y a donc pas de raison d'arrondir le résultat à un nombre entier.
		\end{enumerate}
	\end{enumerate}

\end{document}
